\chapter{The Text Screen}

The display on the \jr\ is managed by the video controller of the TinyVicky chip. The display controller provides four different engines for generating images on the screen:

\begin{itemize}
    \item Text: an old school style text screen where the characters to display are stored in a text matrix, and the shape of those characters comes from font memory. Text mode characters are 8 pixels wide by 8 pixels high.
    \item Bitmap: a simple pixel image engine that can display a bitmap image covering the screen.
    \item Sprite: an engine to display small, movable sprites on the screen.
    \item Tile: an engine to display images on the screen made up of tiles from a tile set.
\end{itemize}

The bitmap, sprite, and tile engines are considered graphics modes. TinyVicky will let you display either text by itself, a mix of the graphics modes by themselves, or text overlaid on top of the graphics modes.

\section*{Text Matrix}

The memory for the characters to display on the screen is the text matrix, which is stored in I/O page 2. When this I/O page is swapped into the CPU address space, it appears at 0x18\_4000. Each byte of memory corresponds to a single character on the screen in left to right, top to bottom order. The byte at 0x18\_4000 (Block \$C2 Offset \$0000) is the upper left corner of the screen, the byte at 0x18\_4001 is the next character to the right, and so on. The number of bytes per line is set by the base resolution of the screen, but is generally 80. When a border is displayed, while that limits the number of characters displayed, the layout in memory remains the same.

The text screen has two core resolutions, tied to the refresh rate of the screen: 80 by 60 at 60 Hz, and 80 by 50 at 70 Hz. Beyond that, the character display may be made double width or double height, or both. This gives the following possible character displays: $80 \times 60$, $40 \times 60$, $80 \times 30$, $40 \times 30$, $80 \times 50$, $40 \times 50$, $80 \times 25$, and $40 \times 25$.

\example{Print an A to the Screen}
\begin{verbatim}
    lda $0001       ; Save the current MMU setting
    pha

    lda #$02        ; Swap I/O Page 2 into bank 6
    sta $0001

    lda #'A'        ; Write 'A' to the upper left corner
    sta $C000

    pla             ; Restore the old MMU setting
    sta $0001
\end{verbatim}

Note: this example does not set the font or the color, so depending on how your \jr\ is initialized, you may not see an actual ``A'' on the screen.

\section*{Text Color LUTs}

Characters in TinyVicky text mode have two colors: the foreground and the background. The foreground and background colors are picked for each character out of two different palettes of 16 colors each. The colors in the palettes are picked from the full range of colors \jr\ can produce, which is more than 16 million colors. This is all managed through two color lookup tables (LUTs) provided by TinyVicky: a text foreground color LUT, and a text background color LUT.

The text LUTs are stored in I/O page 0. The foreground LUT starts at 18x3800, and the background LUT starts at 18x3840 (Block \$C1 0x1800 and Block \$C1 0x1840).

Each LUT is a list of 16 entries. Each entry is a set of four bytes: blue, green, red, and alpha (in that order). Each byte indicates how much of that primary color is present as a component of the actual color. The values range from 0 (none) to 255 (as much as possible). Currently, the alpha channel is not used and is there for future expansion.

\begin{table}[ht]
    \begin{center}
        \begin{tabular}{|c|c|c|c|c|c|c|c|} \hline
            Index & R/W & Foreground & Background & 0 & 1 & 2 & 3 \\ \hline\hline
            0 & W & \verb+18x3800+ & \verb+18x3840+ & BLUE\_0 & GREEN\_0 & RED\_0 & X \\ \hline
            1 & W & \verb+18x3804+ & \verb+18x3844+ & BLUE\_1 & GREEN\_1 & RED\_1 & X \\ \hline
            2 & W & \verb+18x3808+ & \verb+18x3848+ & BLUE\_2 & GREEN\_2 & RED\_2 & X \\ \hline
            3 & W & \verb+18x380C+ & \verb+18x384C+ & BLUE\_3 & GREEN\_3 & RED\_3 & X \\ \hline
            4 & W & \verb+18x3810+ & \verb+18x3850+ & BLUE\_4 & GREEN\_4 & RED\_4 & X \\ \hline
            5 & W & \verb+18x3814+ & \verb+18x3854+ & BLUE\_5 & GREEN\_5 & RED\_5 & X \\ \hline
            6 & W & \verb+18x3818+ & \verb+18x3858+ & BLUE\_6 & GREEN\_6 & RED\_6 & X \\ \hline
            7 & W & \verb+18x381C+ & \verb+18x385C+ & BLUE\_7 & GREEN\_7 & RED\_7 & X \\ \hline
            8 & W & \verb+18x3820+ & \verb+18x3860+ & BLUE\_8 & GREEN\_8 & RED\_8 & X \\ \hline
            9 & W & \verb+18x3834+ & \verb+18x3864+ & BLUE\_9 & GREEN\_9 & RED\_9 & X \\ \hline
            10 & W & \verb+18x3828+ & \verb+18x3868+ & BLUE\_10 & GREEN\_10 & RED\_10 & X \\ \hline
            11 & W & \verb+18x382C+ & \verb+18x386C+ & BLUE\_11 & GREEN\_11 & RED\_11 & X \\ \hline
            12 & W & \verb+18x3830+ & \verb+18x3870+ & BLUE\_12 & GREEN\_12 & RED\_12 & X \\ \hline
            13 & W & \verb+18x3834+ & \verb+18x3874+ & BLUE\_13 & GREEN\_13 & RED\_13 & X \\ \hline
            14 & W & \verb+18x3838+ & \verb+18x3878+ & BLUE\_14 & GREEN\_14 & RED\_14 & X \\ \hline
            15 & W & \verb+18x383C+ & \verb+18x387C+ & BLUE\_15 & GREEN\_15 & RED\_15 & X \\ \hline
        \end{tabular}
    \end{center}
    \caption{Text Color Lookup Tables}
    \label{tab:text_luts}
\end{table}

\begin{table}[ht]
    \begin{center}
        \begin{tabular}{|c|c|c|c|c|c|c|c|} \hline
            Index & R/W & Foreground & Background & 0 & 1 & 2 & 3 \\ \hline\hline
            0 & W & \verb+0x1800+ & \verb+0x1840+ & BLUE\_0 & GREEN\_0 & RED\_0 & X \\ \hline
            1 & W & \verb+0x1804+ & \verb+0x1844+ & BLUE\_1 & GREEN\_1 & RED\_1 & X \\ \hline
            2 & W & \verb+0x1808+ & \verb+0x1848+ & BLUE\_2 & GREEN\_2 & RED\_2 & X \\ \hline
            3 & W & \verb+0x180C+ & \verb+0x184C+ & BLUE\_3 & GREEN\_3 & RED\_3 & X \\ \hline
            4 & W & \verb+0x1810+ & \verb+0x1850+ & BLUE\_4 & GREEN\_4 & RED\_4 & X \\ \hline
            5 & W & \verb+0x1814+ & \verb+0x1854+ & BLUE\_5 & GREEN\_5 & RED\_5 & X \\ \hline
            6 & W & \verb+0x1818+ & \verb+0x1858+ & BLUE\_6 & GREEN\_6 & RED\_6 & X \\ \hline
            7 & W & \verb+0x181C+ & \verb+0x185C+ & BLUE\_7 & GREEN\_7 & RED\_7 & X \\ \hline
            8 & W & \verb+0x1820+ & \verb+0x1860+ & BLUE\_8 & GREEN\_8 & RED\_8 & X \\ \hline
            9 & W & \verb+0x1824+ & \verb+0x1864+ & BLUE\_9 & GREEN\_9 & RED\_9 & X \\ \hline
            10 & W & \verb+0x1828+ & \verb+0x1868+ & BLUE\_10 & GREEN\_10 & RED\_10 & X \\ \hline
            11 & W & \verb+0x182C+ & \verb+0x186C+ & BLUE\_11 & GREEN\_11 & RED\_11 & X \\ \hline
            18 & W & \verb+0x1830+ & \verb+0x1870+ & BLUE\_18 & GREEN\_18 & RED\_18 & X \\ \hline
            18 & W & \verb+0x1834+ & \verb+0x1874+ & BLUE\_18 & GREEN\_18 & RED\_18 & X \\ \hline
            14 & W & \verb+0x1838+ & \verb+0x1878+ & BLUE\_14 & GREEN\_14 & RED\_14 & X \\ \hline
            15 & W & \verb+0x183C+ & \verb+0x187C+ & BLUE\_15 & GREEN\_15 & RED\_15 & X \\ \hline
        \end{tabular}
    \end{center}
    \caption{Text Color Lookup Tables Relative to Block \$C1}
    \label{tab:text_luts_relative}
\end{table}

\section*{Color Matrix}

The way that text color is selected for each character is through the color matrix. This section of memory is in I/O page 3 and starts at 0x18\_6000 (Block \$C3 Offset \$0000) when page 3 is swapped into the CPU's address space. The layout is precisely the same as the text matrix ({\it e.g.} the character at 0x18\_4023 in the text matrix has its color information at 0x18\_6023 in the color matrix).

Each byte in the color matrix specifies two colors by providing an index into each of the two text LUTs. The most significant four bits is the number of the foreground color to use. The number of the least significant four bits is the number of the background color to use.

Let's say the color value at 0x18\_6023 is 0x45. This means that the foreground color of the character is color 4 from the text foreground LUT, which starts at 0xFF10 (0xFF00 + 4 * 4), and the background color of the character is 5 from the text background LUT, which starts at 0xFF54 (0xFF40 + 4 * 5). If the bytes at 0xFF10 are 0x00, 0x80, 0x80, that means the foreground will be a medium yellow. If the bytes at 0xFF54 are 0xFF, 0x00, 0x00, that means the background will be blue.

\example{Make That ``A'' Yellow on Blue}
\begin{verbatim}
    lda $0001       ; Save the MMU state
    pha

    stz $0001       ; Switch in I/O Page #0

    stz $D810       ; Set foreground #4 to medium yellow
    lda #$80
    sta $D811
    sta $D812

    lda #$FF        ; Set background #5 to blue
    sta $D854
    stz $D855
    stz $D856

    lda #$03        ; Switch to I/O page #3 (color matrix)
    sta $0001

    lda #$45        ; Color will be foreground=4, background=5
    sta $C000

    pla             ; Restore the MMU state
    sta $0001
\end{verbatim}

\section*{Entering Text Mode}

Whether text mode is being displayed (and in what resolution) is controlled by the VICKY Master Control Registers (see table~\ref{tab:vky_master_ctrl_reg}). For now, we're going to ignore most of the bits, which are used by other display modes. For text mode, we really only care about the TEXT bit, which needs to be set to turn on the text display. The resolution is controlled by DBL\_Y, DBL\_X, and CLK\_70. If we set 0xFFC0 to 0x01 and 0xFFC1 to 0x00, that will put us into text mode at $80 \times 60$.

\begin{table}[ht]
    \begin{center}
        \begin{tabular}{|c|c|c|c|c|c|c|c|c|c|} \hline
            Address & R/W & 7 & 6 & 5 & 4 & 3 & 2 & 1 & 0 \\ \hline\hline
            \verb+0xFFC0+ & R/W & DIS\_VID & GAMMA & SPRITE & TILE & BITMAP & GRAPH & OVRLY & TEXT \\ \hline
            \verb+0xFFC1+ & R/W & SHWBG & MTXT\_EN & FON\_SET & FON\_OVLY & MON\_SLP & DBL\_Y & DBL\_X & CLK\_70 \\ \hline
        \end{tabular}
    \end{center}
    \caption{VICKY Master Control Registers}
    \label{tab:vky_master_ctrl_reg}
\end{table}

\begin{description}
    \item[TEXT] if set (1), text mode display is enabled

    \item[OVRLY] if set, text will be overlaid on graphics

    \item[GRAPH] if set, one or more of the graphics modes may be used

    \item[BITMAP] if set (and GRAPHICS is set), bitmap graphics may be displayed

    \item[TILE] if set (and GRAPHICS is set), tile graphics may be displayed

    \item[SPRITE] if set (and GRAPHICS is set), sprite graphics may be displayed

    \item[GAMMA] if set, gamma correction is enabled

    \item[DIS\_VID] if set, disable video scanning, giving 100\% bandwitth to CPU      

    \item[CLK\_70] if set, the video refresh will be set to 70 Hz mode (640x400 text resolution, 320x200 graphics). If clear,
        the video refresh will be set to 60 Hz (640x480 text resolution, 320x240 graphics).

    \item[DBL\_X] if set, text mode characters will be twice as wide

    \item[DBL\_Y] if set, text mode characters will be twice as high
    
    \item[MON\_SLP] if set, the monitor SYNC signal will be turned off, putting the monitor to sleep
    
    \item[FON\_OVLY] if clear (0), only the text foreground color will be displayed when text overlays graphics (all background colors will be completely transparent). If set (1), both foreground and background colors will be displayed, except that background color 0 will be transparent. 
    
    \item[FON\_SET] if set (1), the text font displayed will be font set 1. If clear (0), the text font displayed will be font set 0.

    \item[MTXT\_EN] if set, enable Memtext mode (prioritized over block text)

    \item[SHWBG] if set, show background in non-overlay mode in MemText mode      
\end{description}

\section*{Text Fonts}
\label{sec:text_font}

Character shapes (or ``glyphs,'' if you prefer) are defined in font memory, which is in I/O page 1 and starts at 0x18\_2000 (Block \$C1 Offset \$0000). TinyVicky provides for two font sets, and which one is used for text mode is controlled by the FON\_SET bit in the Vicky Master Control Register. Only one will be in use at any given time in normal operation. Font set 0 is from 0x18\_2000 through 0x1827FF. Font set 1 is from 0x18\_2800 through 0x182FFF (Block \$C2 Offset \$0800).

The \jr\ treats each character as a square of pixels, 8 pixels on a side. A pixel may be either in the foreground color for the character or in the background color for the character. The way this is managed is that each character has a sequence of eight bytes in the font memory. Each byte represents a row in the character, and each bit represents a pixel in the row ($\blacksquare$ for foreground, $\square$ for background).

As an example, let's say we wanted to have a fancy ``F'' for character 0:

\begin{table}[ht]
    \begin{center}
        \begin{tabular}{|c|c|c|c|c|c|c|c|c|} \hline
            $\square$ & $\square$ & $\square$ & $\blacksquare$ & $\blacksquare$ & $\blacksquare$ & $\blacksquare$ & $\blacksquare$ & 0x1F \\ \hline
            $\square$ & $\square$ & $\blacksquare$ & $\blacksquare$ & $\square$ & $\square$ & $\square$ & $\square$ & 0x30 \\ \hline
            $\square$ & $\square$ & $\blacksquare$ & $\blacksquare$ & $\square$ & $\square$ & $\square$ & $\square$ & 0x30 \\ \hline
            $\square$ & $\blacksquare$ & $\blacksquare$ & $\blacksquare$ & $\blacksquare$ & $\blacksquare$ & $\square$ & $\square$ & 0x7C \\ \hline
            $\square$ & $\blacksquare$ & $\blacksquare$ & $\square$ & $\square$ & $\square$ & $\square$ & $\square$ & 0x60 \\ \hline
            $\blacksquare$ & $\blacksquare$ & $\square$ & $\square$ & $\square$ & $\square$ & $\square$ & $\square$ & 0xC0 \\ \hline
            $\blacksquare$ & $\blacksquare$ & $\square$ & $\square$ & $\square$ & $\square$ & $\square$ & $\square$ & 0xC0 \\ \hline
        \end{tabular}
    \end{center}
    \caption{A sample character}
    \label{tab:text_font}
\end{table}

The glyph to display would be defined by the eight byte sequence 0x1F, 0x30, 0x30, 0x7C, 0x60, 0xC0, 0xC0. We would store that sequence starting at 0x18\_2000 (0x1F), through 0x18\_2007 (0xC0). After that was set, any time the byte 0x00 is written to screen memory, the glyph ``F'' would be displayed in that position.

\section*{Text Cursor}

\jr\ has a text mode cursor. The text mode cursor is implemented as a character which is displayed in a $(x, y)$ position on the screen, visually replacing the character ordinarily at that position. It may be displayed continuously, or it may flash at one of four rates. When flashing, that position in the text screen will alternate between the text cursor and the character at that position in the text matrix. The color for the text cursor comes from the color for the position on the screen as specified in the color matrix. In other words, the text cursor does not have its own color. The text cursor registers are located in I/O bank 0.

\begin{table}[ht]
    \begin{center}
        \begin{tabular}{|c|c|c|c|c|c|c|c|c|c|c|} \hline
            Address & R/W & Name & 7 & 6 & 5 & 4 & 3 & 2 & 1 & 0 \\\hline\hline
            \verb+0xFFD0+ & R/W & CCR & \multicolumn{4}{|c|}{---} & FLASH\_DIS & \multicolumn{2}{|c|}{RATE} & ENABLE \\ \hline
            \verb+0xFFD2+ & R/W & CCH & \multicolumn{8}{|c|}{Cursor character} \\ \hline
            \verb+0xFFD3+ & R/W & CCO & \multicolumn{8}{|c|}{Cursor Color} \\ \hline            
            \verb+0xFFD4+ & R/W &  \multirow{2}{*}{CURX} & X15 & X14 & X13 & X12 & X11 & X10 & X9 & X8 \\ \cline{1-2}\cline{4-11}
            \verb+0xFFD5+ & R/W &  & X7 & X6 & X5 & X4 & X3 & X2 & X1 & X0 \\ \hline 
            \verb+0xFFC6+ & R/W &  \multirow{2}{*}{CURY} & Y15 & Y14 & Y13 & Y12 & Y11 & Y10 & Y9 & Y8 \\ \cline{1-2}\cline{4-11}
            \verb+0xFFD7+ & R/W &  & Y7 & Y6 & Y5 & Y4 & Y3 & Y2 & Y1 & Y0 \\ \hline
        \end{tabular}
    \end{center}
    \caption{Text Cursor Registers}
    \label{tab:txt_crsr_reg}
\end{table}

\begin{description}
    \item[ENABLE] if this flag is set (1), the cursor is enabled

    \item[FLASH\_DIS] if this flag is set (1), the cursor will not flash. If clear (0), it will flash.

    \item[RATE] these two bits set the rate at which the cursor flashes (see table~\ref{tab:txt_crsr_rates})

    \item[CCH] the character code for the cursor character to display

    \item[CURX] the column number (16-bit) for the cursor

    \item[CURY] the row number (16-bit) for the cursor
\end{description}

\begin{table}[ht]
    \begin{center}
        \begin{tabular}{|c|c|c|} \hline
            RATE1 & RATE0 & Rate \\\hline\hline
            0 & 0 & 1s \\ \hline
            0 & 1 & $1/2$s \\ \hline
            1 & 0 & $1/4$s \\ \hline
            1 & 1 & $1/5$s \\ \hline
        \end{tabular}
    \end{center}
    \caption{Text Cursor Flash Rates}
    \label{tab:txt_crsr_rates}
\end{table}

\section*{MemText}


MemText mode features include: four 8x8 text fonts or two 8x16 text fonts, two foreground CLUTs and two background CLUTs of 256 colors each, and 256 colors per character.  MemText text mode keeps the text and color data in RAM, which allows the use of DMA to copy data and to switch text screen by changing the pointer values to a different value in memory.

\begin{description}
  \item[MemText Registers] located in block \$C0
\end{description}
\begin{table}[ht]
    \begin{center}
        \begin{tabular}{|c|c|c|c|c|c|c|c|c|c|c|} \hline
            Address & R/W & Name & 7 & 6 & 5 & 4 & 3 & 2 & 1 & 0 \\\hline\hline
            \verb+18x1200+ & R/W & MTCR & \multicolumn{6}{|c|}{---} & FSIZE & MTXT\_EN \\ \hline
            \verb+18x1201+ & R/W & CCR & X & CSIZE & TF\_1 & TF\_0 & F\_ON & CF\_1 & CF\_0 & CURS\_EN  \\ \hline
            \verb+18x1202+ & R/W & CURX & X7 & X6 & X5 & X4 & X3 & X2 & X1 & X0 \\ \hline
            \verb+18x1203+ & R/W & CURY & Y7 & Y6 & Y5 & Y4 & Y3 & Y2 & Y1 & Y0 \\ \hline
            \verb+18x1204+ & R/W & N/A & \multicolumn{8}{|c|}{NOT USED} \\ \hline            
            \verb+18x1205+ & R/W & TXTH & T23 & T22 & T21 & T20 & T19 & T18 & T17 & T16 \\ \hline
            \verb+18x1206+ & R/W & TXTM & T15 & T14 & T13 & T12 & T11 & T10 & T9 & T8 \\ \hline
            \verb+18x1207+ & R/W & TXTK & T7 & T6 & T5 & T4 & T3 & T2 & T1 & T0 \\ \hline
            \verb+18x1208+ & R/W & N/A & \multicolumn{8}{|c|}{NOT USED} \\ \hline
            \verb+18x1209+ & R/W & COLH & C23 & C22 & C21 & C20 & C19 & C18 & C17 & C16 \\ \hline
            \verb+18x120A+ & R/W & COLM & C15 & C14 & C13 & C12 & C11 & C10 & C9 & C8 \\ \hline
            \verb+18x120B+ & R/W & COLL & C7 & C6 & C5 & C4 & C3 & C2 & C1 & C0 \\ \hline
            \verb+18x120C+ & R/W & PRE & \multicolumn{8}{|c|}{Pre-Charge (Reserved)} \\ \hline
            \verb+18x120D+ & R/W & CCB & \multicolumn{8}{|c|}{Cursor Color B} \\ \hline
            \verb+18x120E+ & R/W & CCG & \multicolumn{8}{|c|}{Cursor Color G} \\ \hline
            \verb+18x120F+ & R/W & CCR & \multicolumn{8}{|c|}{Cursor Color R} \\ \hline
            \verb+18x1210+ & R/W & CG8 & \multicolumn{8}{|c|}{Cursor Graphic x8} \\ \hline
            \verb+18x1211+ & R/W & CG8 & \multicolumn{8}{|c|}{Cursor Graphic x8} \\ \hline
            \verb+18x1212+ & R/W & CG8 & \multicolumn{8}{|c|}{Cursor Graphic x8} \\ \hline
            \verb+18x1213+ & R/W & CG8 & \multicolumn{8}{|c|}{Cursor Graphic x8} \\ \hline
            \verb+18x1214+ & R/W & CG8 & \multicolumn{8}{|c|}{Cursor Graphic x8} \\ \hline
            \verb+18x1215+ & R/W & CG8 & \multicolumn{8}{|c|}{Cursor Graphic x8} \\ \hline
            \verb+18x1216+ & R/W & CG8 & \multicolumn{8}{|c|}{Cursor Graphic x8} \\ \hline
            \verb+18x1217+ & R/W & CG8 & \multicolumn{8}{|c|}{Cursor Graphic x8} \\ \hline
            \verb+18x1218+ & R/W & CG16 & \multicolumn{8}{|c|}{Cursor Graphic x16} \\ \hline
            \verb+18x1219+ & R/W & CG16 & \multicolumn{8}{|c|}{Cursor Graphic x16} \\ \hline
            \verb+18x121A+ & R/W & CG16 & \multicolumn{8}{|c|}{Cursor Graphic x16} \\ \hline
            \verb+18x121B+ & R/W & CG16 & \multicolumn{8}{|c|}{Cursor Graphic x16} \\ \hline
            \verb+18x121C+ & R/W & CG16 & \multicolumn{8}{|c|}{Cursor Graphic x16} \\ \hline
            \verb+18x121D+ & R/W & CG16 & \multicolumn{8}{|c|}{Cursor Graphic x16} \\ \hline
            \verb+18x121E+ & R/W & CG16 & \multicolumn{8}{|c|}{Cursor Graphic x16} \\ \hline
            \verb+18x121F+ & R/W & CG16 & \multicolumn{8}{|c|}{Cursor Graphic x16} \\ \hline
        \end{tabular}
    \end{center}
    \caption{MemText Registers}
    \label{tab:memtext_reg}
\end{table}

Screen resolutions are 80x60 for 8x8 fonts and 80x30 for 8x16 fonts.  The double pixel registers do not affect MemText mode.

\begin{description}
  \begin{samepage}
    \item[MTCR] Memory Text Control Register
      \begin{description}
      \item[MTXT\_EN] 1=enabled, 0=disabled
      \item[FSIZE] Font Size: 0=8x8, 1=8x16 (State-Machine)      
      \end{description}
  \end{samepage}
    \item[CCR] Cursor Control Register
      \begin{description}
      \item[CURS\_EN] 1=enabled, 0=disabled
      \item[CF\_0] Cursor Flashing Rate 0
      \item[CF\_1] Cursor Flashing Rate 1
      \item[F\_ON] Cursor Flashing on/off
      \item[TF\_0] Text Flashing Rate 0 (triggered by attributes)
      \item[TF\_1] Text Flashing Rate 1 (triggered by attributes)
      \item[CSIZE]  Cursor Size: 0=8x8, 1=8x16
        \begin{table}[ht]
    \begin{center}
        \begin{tabular}{|c|c|c|} \hline
            RATE1 & RATE0 & Rate \\\hline\hline
            0 & 0 & 1bps \\ \hline
            0 & 1 & 2bps \\ \hline
            1 & 0 & 3bps \\ \hline
            1 & 1 & 5bps \\ \hline
        \end{tabular}
    \end{center}
    \caption{Text Cursor Flash Rates}
    \label{tab:mtext_crsr_rates}
\end{table}
      \end{description}

    \item[CURX] Cursor X Register

    \item[CURY] Cursor Y Register

    \item[TXTH] Text Address Pointer H
      
    \item[TXTM] Text Address Pointer M

    \item[TXTL] Text Address Pointer L

    \item[COLH] Color Address Pointer H

    \item[COLM] Color Address Pointer M

    \item[COLL] Color Address Pointer L

    \item[PRE] Pre-Charge (Reserved - defaulted to \$0A)

    \item[CCB] Cursor Color B

    \item[CCG] Cursor Color R

    \item[CCR] Cursor color R

    \item[CG8] Cursor Graphics x8

    \item[CG16] Cursor Graphics x16    
\end{description}

\section*{MemText Foreground/Background LUTs}

The foreground and background LUTs consist of 256 colors in BGRA format.  Each LUT is 1K.

\begin{table}[ht]
    \begin{center}
        \begin{tabular}{|c|c|c|c|c|} \hline
          Block & Offset & Address Start & Address End & Description \\\hline\hline
          \$C5 & \verb+0x0000+ & \verb+18xA000+ & \verb+18xA3FF+ & MEMTEXT FOREGROUND - LUT0 \\ \hline
          \$C5 & \verb+0x0400+ & \verb+18xA400+ & \verb+18xA7FF+ & MEMTEXT FOREGROUND - LUT0 \\ \hline
          \$C5 & \verb+0x0800+ & \verb+18xA800+ & \verb+18xABFF+ & MEMTEXT FOREGROUND - LUT0 \\ \hline
          \$C5 & \verb+0x0C00+ & \verb+18xAC00+ & \verb+18xAFFF+ & MEMTEXT FOREGROUND - LUT0 \\ \hline
       \end{tabular}
    \end{center}
    \caption{MemText Foreground/Background LUTS}
    \label{tab:memtext_fgbglut}
\end{table}

\section*{MemText Fonts}

\begin{table}[ht]
    \begin{center}
        \begin{tabular}{|c|c|c|c|c|} \hline
          Block & Offset & Address Start & Address End & Description \\\hline\hline
          \$C6 & \verb+0x0000+ & \verb+18xC000+ & \verb+18xDFFF+ & MEMTEXT FONT (2x8x16),(4x8x8) \\ \hline
        \end{tabular}
    \end{center}
    \caption{MemText Registers}
    \label{tab:memtext_fonts}
\end{table}

\section*{MemText Attributes}

The text memory for MemText is 2 bytes per character.  The first byte is the character, the second byte contains the attributes.

\begin{table}[ht]
    \begin{center}
        \begin{tabular}{|c|c|c|c|c|c|c|c|c|} \hline
          Character &  \multicolumn{8}{|c|}{Attributes} \\\hline\hline 
          7:0 & 7 & 6 & 5 & 4 & 3 & 2 & 1 & 0 \\ \hline 
          Character & X & X & FL\_EN & INV & BGLUT & FGLUT & FSEL1 & FSEL0 \\ \hline 
          
        \end{tabular}
    \end{center}
    \caption{MemText Registers}
    \label{tab:memtext_characters}
\end{table}

\begin{description}
    \item[FSEL1 and FSEL0] selects the font bank.  If 8x16 font, then only FSEL0 is used to select one of two font banks.  For 8x8, FSEL0 and FSEL1 are used to select one of four font banks

    \item[FGLUT] select foreground LUT0 or LUT1

    \item[INV] inverse character

    \item[FL\_EN] 1=enable flashing, 0=disable flashing

\end{description}
